\section{Structural derivations}
\label{sec:structural}

\subsection{Delay-chain moments and pure-delay limit}
Let $r=n/\theta$ be the delay-state rate. The normalized chain kernel for
$t\geq0$ is
\[
 h_n(t)=\frac{r^n t^{n-1}e^{-rt}}{(n-1)!}.
\]
Its Laplace transform is $(1+s/r)^{-n}$. Differentiating that transform at zero
gives mean $n/r=\theta$ and variance $n/r^2=\theta^2/n$. Consequently, the
normalized complete thermal response has impulse-time mean $\tau+\theta$ and
variance $\tau^2+\theta^2/n$. These are properties of the model kernel;
they are not estimates of a measured residence-time distribution.
For fixed $s$ in the domain of convergence,
\[
 \log(1+\theta s/n)^{-n}=-n\log(1+\theta s/n)\longrightarrow-\theta s.
\]
This calculation supplies the pure-delay limit of Eq.~\eqref{eq:transfer}.

\subsection{An explicit physical-parameter equivalence class}
Define the two positive pole time constants $t_{\mathrm f}<t_{\mathrm s}$ by
$a=t_{\mathrm f}+t_{\mathrm s}$ and $b=t_{\mathrm f}t_{\mathrm s}$. For fixed
positive transfer coefficients $(K,a,b)$, choose a family coordinate
$r=\CpS/\Ub$ between these pole time constants. The physical parameters
\begin{equation}
 \Ua=\frac{c}{K},\qquad
 \CpH=\frac{b\Ua}{r},\qquad
 \CpS=\Ua\left(a-r-\frac{b}{r}\right),\qquad
 \Ub=\frac{\CpS}{r}
 \label{eq:equivalent}
\end{equation}
produce exactly the same transfer function for every
$r\in(t_{\mathrm f},t_{\mathrm s})$. Positivity follows from
$a-r-b/r=-(r-t_{\mathrm f})(r-t_{\mathrm s})/r>0$. Substitution into
Eq.~\eqref{eq:original-transfer} recovers $K$, $a$, and $b$ exactly.
The conductance $\Ua$ is identifiable through $K$; the other three parameters
vary along this equivalence class. No input waveform can distinguish members
of this class through this output under the stated initial conditions.

Positivity constrains this continuum to finite ranges in some coordinates.
For example, its heat capacities obey
\[
 \Ua t_{\mathrm f}<\CpH<\Ua t_{\mathrm s},\qquad
 0<\CpS\leq\Ua\left(a-2\sqrt b\right).
\]
The second upper bound occurs at $r=\sqrt b$.
Thus bounded physical ranges or bounded exploration grids do not establish
structural identifiability. Structural uniqueness must be checked separately
from the numerical width of a confidence set.

\subsection{Constant input with equilibrated delay states}
If $u(t)=u_0$ and $z_i(0)=u_0$, every delay-state derivative in
Eq.~\eqref{eq:chain} vanishes. The predicted temperature then depends on
$K$ and $\tau$ alone, and its derivative with respect to $\theta$ is zero at
every observation time. A constant prehistory gives the same conclusion for
Eq.~\eqref{eq:puredelay}. This experiment cannot estimate the delay even when
the thermal transient is well measured. Cold delay states under the same
post-start input retain a transient that depends on $\theta$.

\subsection{Coordinate changes and information}
For positive parameters, write $\bm\eta=(\log K,\log\tau,\log\theta)^\top$
and $\bm\xi=(\log\Ua,\log\Cp,\log\theta)^\top$. The Jacobian of the map
from $\bm\eta$ to $\bm\xi$ is the constant matrix
\[
 M=\begin{pmatrix}-1&0&0\\-1&1&0\\0&0&1\end{pmatrix},\qquad M^{-1}=M.
\]
For a nonsingular covariance $C_{\eta}$, the transformed local covariance is
$C_{\xi}=M C_{\eta}M^\top$. Sensitivity rank and prediction-space uncertainty
are invariant under this invertible transformation. Numerical eigenvalues,
eigenvectors, and condition numbers generally change with the coordinate metric.
The matrix satisfies $|\det M|=1$, so log-coordinate information determinants
at the same model point are equal in this particular transformation.
